%! Author = murfe
%! Date = 29.06.2026
\chapter{Anforderungen}
\label{ch:anforderungen}
%TODO: Einleitung
\section{Abgrenzung des betrachteten Fehlermodells}
\label{sec:abgrenzungFehlermodell}
Verteilte Softwaresysteme können durch eine Vielzahl unterschiedlicher Fehler beeinträchtigt werden.
Hierzu zählen unter anderem Hardwaredefekte, Softwarefehler, Netzwerkausfälle, Bedienfehler oder Cyberangriffe.
Wie bereits in \cref{sec:fehlertoleranz} erläutert wurde, kann ein System jedoch nicht sämtliche denkbaren Fehler verhindern oder behandeln.
Der Entwurf fehlertoleranter Systeme erfordert daher eine bewusste Abwägung, welche Fehlerarten berücksichtigt und welche Risiken akzeptiert werden.
Diese Entscheidung orientiert sich zum am Anwendungskontext und den Anforderungen an das System und zum anderen an wirtschaftlichen Randbedingungen
wie Entwicklungs- und Betriebsaufwand \parencite[Kap.~1 Abschnitt \glqq Reliability\grqq]{kleppmann_designing_2017}.

Die im Folgenden dargelegte Teilmenge betrachteter Fehler ist eine Auswahl der von
\textcite{kleppmann_designing_2017,newman_building_2026,bass_software_2021} beschriebenen typisch vorkommenden Fehlern in verteilten Systemen.
Ziel dieser Arbeit ist der Entwurf einer fehlertoleranten Architektur für ein datenbankgestütztes Background Job-Processing-System.
Betrachtet werden daher ausschließlich Fehlerklassen, die während des regulären Betriebs auftreten können und deren Auswirkungen durch Mechanismen der
Anwendungsarchitektur beherrscht werden können.
Fehler, deren Behandlung zusätzliche Infrastruktur oder organisatorische Maßnahmen wie Backups, Replikation oder physische Redundanz erfordern,
liegen außerhalb des Betrachtungsumfangs dieser Arbeit.

Die Abgrenzung der betrachteten Fehler ist in \cref{tab:fehlermodell} zusammengefasst.
\begin{table}
    \centering
    \begin{tabularx}{\textwidth}{lcX}
        \toprule
        \textbf{Fehlerklasse} & \textbf{Relevant} & \textbf{Begründung} \\
        \midrule
        Komponenten und Infrastrukturausfälle & \checkmark & Kurze Ausfälle können durch Software behandelt werden\\
        Doppelte Requests & \checkmark & Mehrfachübermittlung desselben Jobs\\
        Doppelte Ausführung & \checkmark & Wiederaufnahme begonnener Jobs \\
        Konkurrierender Zugriff & \checkmark & Parallele Worker können Race Conditions verursachen \\
        Physischer Datenbankverlust & \ding{55} & Erfordert Backups \& Redundanz \\
        Totalausfall Infrastruktur & \ding{55} & Erfordern Disaster-Recovery Maßnahmen \\
        Cyberangriffe & \ding{55} & Qualitätsmerkmal Sicherheit\\
        Byzantinische Fehler & \ding{55} & Annahme: Vertrauensvolle Komponenten\\
        \bottomrule
    \end{tabularx}
    \caption{}
    \label{tab:fehlermodell}
\end{table}
Innerhalb des Systementwurfs werden im Wesentlichen vier Kategorien von Fehlern behandelt.
An erster Stelle stehen Ausfälle einzelner Systemkomponenten sowie transiente Infrastruktur- und Kommunikationsstörungen.
Hierzu zählen insbesondere das unerwartete Beenden eines Worker-Prozesses, Container-Neustarts, JVM-Abstürze, durch das Betriebssystem terminierte Prozesse,
temporäre Netzwerkunterbrechungen oder kurzfristig nicht erreichbare Datenbanken.
Dadurch können zwei Arten von Fehlern passieren.
Zum einen besteht die Gefahr, dass eine laufende Verarbeitung unterbrochen wird oder vorübergehend nicht fortgesetzt werden kann.
Aufgrund der Charakteristik transienter Fehler wird dabei grundsätzlich davon ausgegangen, dass die betroffene Komponente oder
Ressource nach kurzer Zeit wieder verfügbar ist.
Zum anderen kann es sein, dass eine Nachricht bzw. ein Job verloren geht.
Durch diese Fehler gerät das Gesamtsystem in einen unklaren Zustand darüber, wie weit die Verarbeitung bereits fortgeschritten war
bzw. ob der Job persistiert wurde.
Da diese Störungen in verteilten Umgebungen regelmäßig auftreten können \parencite[Kap.~1 Abschnitt \glqq Failure Is Inevitable\grqq]{newman_building_2026},
muss die Systemarchitektur in der Lage sein, diese autonom zu überbrücken.

Als direkte Folge solcher Infrastrukturprobleme und der darauf reagierenden Wiederholungsmechanismen der Clients entstehen doppelte Requests.
Ein Client sendet eine Anfrage erneut, weil eine Bestätigung aufgrund eines Timeouts ausblieb.
Die Architektur muss in diesem Szenario durch Duplikaterkennung verhindern, dass identische Jobs mehrfach im Job-Store persistiert werden.

Eng damit verknüpft ist die doppelte Ausführung von Jobs.
Bei der Wiederaufnahme eines zuvor abgebrochenen Jobs kann nach einem Fehler oft nicht eindeutig festgestellt werden, ob die vorherige Verarbeitung
bereits Seiteneffekte erzielt hat.
Die Architektur muss daher zwingend sicherstellen, dass wiederholte Ausführungen keine inkonsistenten Ergebnisse verursachen.

Des Weiteren werden konkurrierende Zugriffe betrachtet.
Da die Parallelverarbeitung mittels mehrerer, autonom agierender Worker ein wesentliches Merkmal moderner Background Job-Processing-Systeme darstellt,
entstehen unweigerlich gleichzeitige Zugriffe auf geteilte Datenstrukturen wie den zentralen Job-Store.
Dies ist im klassischen Sinne zwar nicht direkt ein Fehler, jedoch können parallele Zugriffe ohne entsprechende Synchronisationsmechanismen
potenziell zu Race Conditions und inkonsistenten Zuständen in der Persistenzschicht führen, weshalb deren Verhinderung wesentlicher Bestandteil des Entwurfs ist.

Dem gegenüber stehen Fehler, die bewusst außerhalb des Betrachtungsumfangs dieser Arbeit liegen.
Ein physischer Datenbankverlust sowie ein Totalausfall der Infrastruktur können nicht durch die Softwarearchitektur des Job-Processing-Systems gelöst werden.
Ein Verlust der Persistenzschicht oder ein Ausfall des gesamten Rechenzentrums infolge von großflächigen Stromausfällen oder Bränden am Serverstandort
erfordert stattdessen infrastrukturelle Maßnahmen wie geografische Redundanz, kontinuierliche Backups und umfassende Disaster-Recovery-Konzepte.

Ebenso werden Cyberangriffe im Rahmen dieser Arbeit ausgeschlossen.
Obwohl Angriffe wie SQL-Injections oder Denial-of-Service-Szenarien direkte Auswirkungen auf die Zuverlässigkeit eines Systems haben können, betreffen sie
primär das Qualitätsmerkmal der Sicherheit.
Ihre Abwehr erfordert spezialisierte Mechanismen wie strikte Eingabevalidierung, Ratenbegrenzungen, Authentifizierung und Angriffserkennungssysteme.
Der Fokus dieser Arbeit liegt stattdessen auf der Gewährleistung von Fehlertoleranz bei regulärem Systembetrieb.

Schließlich werden auch byzantinische Fehler ausgeschlossen.
Dies sind Fehler, bei denen kompromittierte oder defekte Komponenten willkürlich falsche oder manipulierte Informationen erzeugen und teilen
\parencite[Kap.~8 Abschnitt \glqq Byzantine Faults\grqq]{kleppmann_designing_2017}.
Die vorliegende Arbeit geht dahingehend von der Annahme ehrlicher Komponenten aus.
Dies bedeutet, dass Komponenten zwar abstürzen oder vorübergehend nicht erreichbar sein können, im aktiven Zustand jedoch nie absichtlich
falsche Informationen verbreiten.

Aus der gewählten Abgrenzung ergibt sich, dass der Schwerpunkt dieser Arbeit auf der Behandlung von Teilausfällen, konkurrierender Verarbeitung und
transienten Infrastrukturfehlern liegt.
Diese Fehlerszenarien besitzen unmittelbaren Einfluss auf die Zuverlässigkeit datenbankgestützter Background Job-Processing-Systeme und bilden daher die
Grundlage für die im Folgenden formulierten funktionalen Anforderungen, Qualitätsanforderungen und Systemzusicherungen.

\section{Funktionale Anforderungen}
\label{sec:funktionaleAnf}
FA1
Das System muss Jobs über REST entgegennehmen.
Begründung:
Da das Zielsystem als Backend-Service konzipiert wird, erfolgt die Kommunikation über HTTP.

FA2
Das System muss Jobs persistent speichern.
Begründung:
Ansonsten würden Jobs bei Prozessabstürzen verloren gehen.

FA3
Das System muss mehrere Worker unterstützen.
Begründung:
Parallelität erhöht den Durchsatz.

FA4
Das System muss Jobs nach Worker-Ausfällen erneut ausführen können.
Begründung:
Worker können jederzeit abstürzen.

FA5
Das System muss den Zustand eines Jobs verwalten.
Begründung:
Recovery setzt Kenntnis des aktuellen Zustands voraus.


\section{Qualitätsanforderungen}
\label{sec:qualitAnf}
QA1
Konsistenz
Fachliche Daten dürfen trotz Fehlern niemals inkonsistent werden.
Begründung:
Grundvoraussetzung für zuverlässige Verarbeitung.

QA2
Fehlertoleranz
Das System soll transiente Fehler automatisch behandeln.
Begründung:
Netzwerk- und Infrastrukturfehler sind normal.

QA3
Recoverability
Jobs dürfen durch Worker-Abstürze nicht verloren gehen.
Begründung:
Kernziel der Arbeit.

QA4
Skalierbarkeit
Mehrere Worker sollen parallel arbeiten können.
Begründung:
Erhöhung des Durchsatzes.

QA5
Nachvollziehbarkeit
Der Verarbeitungszustand eines Jobs muss jederzeit bestimmbar sein.
Begründung:
Recovery und Monitoring.


\section{Zusicherungen}
\label{sec:zusicherungen}

Z1
Ein Job befindet sich zu jedem Zeitpunkt in genau einem gültigen Zustand.

Z2
Ein Job wird niemals dauerhaft verloren.

Z3
Ein Job wird niemals gleichzeitig von mehreren Workern verarbeitet.

Z4
Nach erfolgreichem Abschluss ist das fachliche Ergebnis genau einmal vorhanden.
(oder At-Least-Once + Idempotenz, je nach Formulierung)

Z5
Nach einem Worker-Absturz kann der Job erneut verarbeitet werden.

Z6
Ein erfolgreich abgeschlossener Job wird nicht erneut gestartet.

Z7
Zwischenzustände bleiben nach Rollback nicht dauerhaft bestehen.

Sie ergeben sich unmittelbar aus den zuvor definierten funktionalen Anforderungen und den betrachteten Fehlerszenarien.


Z.B.:
Für die Evaluation werden folgende Zusicherungen definiert, die das System zu jedem Zeitpunkt einhalten muss:
Z1: Kein Job erreicht mehr als einmal den Status SUCCEEDED.
Z2: Zu einem gegebenen Auftrag (Idempotency-Key) wird höchstens ein Report erzeugt.
Z3: Statusübergänge folgen ausschließlich dem definierten Zustandsmodell. Ungültige Übergänge (z. B. direkt von QUEUED zu SUCCEEDED) sind ausgeschlossen.
Z4: Kein Job verbleibt dauerhaft in einem Zwischenzustand. Jeder Job erreicht innerhalb einer definierten Frist einen Endzustand (SUCCEEDED oder FAILED).


Begründen, warum ich mich für diese Frameworks entschieden habe, die vorkommen.
--> Jede (auch implizite) Entscheidung begründen
Alles Begründen!!! Wieso ist das überhaupt relevant? Welches Problem löse ich, wieso habe ich das so gemacht?
Wieso genau diese Fehlerszenarien? Begründen wieso nicht andere oder mehr?
Gibt es formale Modelle die die Zusicherungen zu jedem Zeitpunkt checken?